\subsection{Tame height two at minimal cyclotomic ramification}

\begin{lemma}[A rational second-level point forces a wild absolute index]
\label{lem:B-RD49-point}
\status{Proved; complete root and obstruction/crystal proof, original-source, and exact-code audits passed.}
Let $R$ be a complete discrete valuation ring of mixed characteristic
$(0,p)$ with algebraically closed residue field, and normalize its
valuation by $v(R^\times)=0$ and $v(\pi)=1$. Put $a=v(p)$.
Let $F/R$ be a one-dimensional formal group whose special fibre has
height two. If
\[
             a\leq(p-1)(p^2-1),\qquad p\nmid a,
                                                        \tag{B-RD49.1}
\]
then $F(\mathfrak m_R)$ has no point of exact order $p^2$.
\end{lemma}
\begin{proof}
Choose an identification of the special fibre with the height-two
Honda group, using \cite[Theorem~IV, pp.270--271]{Lazard1955Formal}.
The $p$-typical universal deformation gives an integral coordinate
and $u\in\mathfrak m_R$ in which
\[
 [p](T)=pT+c_pT^p+\sum_{j>p}c_jT^j,\qquad
 c_p=(1-p^{p-1})u,\qquad
 c_j\in(p,u)R\ (p<j<p^2),\qquad c_{p^2}\in R^\times.
                                                        \tag{B-RD49.2}
\]
All coefficients are integral. Here the precise source is
\cite[Proposition~5.7 and its displayed congruence; equations
(12.3)--(12.5) and Proposition~12.10 with proof]{HopkinsGross1994Bundles}.
The latter proposition explicitly allows a complete local
$\mathbf Z_p$-algebra whose residue field extends $\mathbf F_p$;
there is no finite-residue-field restriction on $R$.
For clarity, the logarithm starts
$f(T)=T+(u/p)T^p+O(T^{p^2})$. Comparing
$f([p]T)=pf(T)$ gives the zero coefficients in degrees $2$ through
$p-1$ and the displayed exact value of $c_p$.
Proposition~5.7, with its second coefficient specialized to $1$,
gives the remaining coefficient ideal and the unit coefficient.
An integral invertible change of parameter preserves positive
valuations, so it does not change the question about rational
formal points.

Every nonzero point $z$ of order $p$ satisfies
\[
                           v(z)\leq a/(p-1).            \tag{B-RD49.3}
\]
Indeed, if this inequality failed, $pz$ would have strictly smaller
valuation than every other term of (B-RD49.2), an impossibility.
Now suppose $x$ has exact order $p^2$, and write
$A=v(x)\geq1$ and $Z=v([p]x)$. By (B-RD49.1) and (B-RD49.3),
\[
              Z\leq a/(p-1)\leq p^2-1,
              \qquad a+A>Z,\qquad p^2A>Z.             \tag{B-RD49.4}
\]
Put $b=v(u)$, allowing $b=+\infty$ when $u=0$.
If $b\geq a$, the linear term and every term of degree less than
$p^2$ have valuation at least $a+A$, while all later terms have
valuation at least $p^2A$. This contradicts (B-RD49.4).
Thus $1\leq b<a$. Every term strictly between degrees $p$ and
$p^2$ has valuation at least $b+jA>b+pA$; every term beyond degree
$p^2$ has valuation greater than $p^2A$. Therefore (B-RD49.4)
forces the unique contributing minimum to be the degree-$p$ term:
\[
                              Z=b+pA.                  \tag{B-RD49.5}
\]
At the point $[p]x$ of order $p$, one has
$(p^2-p)Z>b$, so the degree-$p^2$ term and all subsequent terms
have greater valuation than the degree-$p$ term. The intermediate
terms also have greater valuation. Cancellation consequently
requires the linear and degree-$p$ terms to have equal valuation:
\[
            a+Z=b+pZ,\qquad
            a=b+(p-1)Z=p\bigl(b+(p-1)A\bigr).
                                                        \tag{B-RD49.6}
\]
This contradicts $p\nmid a$.
\end{proof}

\begin{lemma}[Tame formal inertia for $p\geq5$]
\label{lem:B-RD49-inertia}
\status{Proved; complete root and obstruction/crystal proof, original-source, and exact-code audits passed.}
Let $p\geq5$ and let $K$ be complete discretely valued of mixed
characteristic $(0,p)$ with algebraically closed residue field.
A height-two BT group which acquires one-dimensional formal good
reduction has a faithful good-reduction extension which is tame
cyclic of degree dividing $p^2-1$. Two such groups have simultaneous
good formal reduction over a tame cyclic extension of degree
$m\mid p^2-1$.
\end{lemma}
\begin{proof}
The endomorphism algebra of the special height-two Honda group is
\[
 D=\mathbf Q_{p^2}\langle\Pi\rangle,
 \qquad \Pi^2=p,\qquad \Pi a=\sigma(a)\Pi,
 \quad a\in\mathbf Q_{p^2}.
                                                        \tag{B-RD49.7}
\]
Here $\sigma$ is the nontrivial unramified field automorphism. This is the same direct Dieudonne calculation
as (B-RD31.1), with $3$ replaced by $p$.
The maximal order has residue field $\mathbf F_{p^2}$, and the
kernel of reduction on its units is pro-$p$.
Every element of $D$ satisfies its reduced quadratic polynomial.
An element of order $p$ would generate
$\mathbf Q_p(\zeta_p)$ of degree $p-1>2$, so cannot occur.
Cauchy's theorem therefore makes every finite subgroup of the
unit group prime to $p$, and reduction embeds it in the cyclic
group $\mathbf F_{p^2}^{\times}$.

The kernel of the finite special inertia action can be removed
from a good-reduction extension by the crystalline descent and
actual-lattice argument of Lemma~\ref{lem:B-RD17-degree}.
Thus the faithful quotient of order dividing $p^2-1$ gives good
reduction for the original lattice, and its special fibre is still
of slope $1/2$. Over this strictly henselian complete base such an
extension is tame cyclic. The tame inertia group is procyclic,
so the common extension of two of them has degree equal to the
least common multiple of their degrees, again dividing $p^2-1$.
\end{proof}

\begin{proposition}[The mixed case is already detected at $p^2$]
\label{prop:B-RD49-mixed}
\status{Proved; complete root and obstruction/crystal proof, original-source, and exact-code audits passed.}
Let $p\geq5$, let $K$ be a complete discretely valued field of mixed
characteristic $(0,p)$ with perfect residue field, and suppose
$e_K=p-1$ and $\mu_p\subset K$.
Let $M$ be a height-two group which acquires supersingular good
reduction. Let $N$ admit a filtration by two potentially good
height-one groups whose dimensions sum to one. Then
\[
                            M[p^2]\not\simeq N[p^2].    \tag{B-RD49.8}
\]
\end{proposition}
\begin{proof}
Pass to the completed maximal unramified extension. Its absolute
ramification is unchanged, its residue field is algebraically
closed, and the torsion comparison would remain an isomorphism.
A potentially good height-one character is
$\chi_p^b\epsilon$, where $b\in\{0,1\}$ and
$\epsilon$ takes values in $\mu_{p-1}\subset\mathbf Z_p^\times$.
The determinant of $T_pM$ is $\chi_p\epsilon_0$ with
$\epsilon_0$ another such finite character. These are the same
height-one and determinant arguments used in
Propositions~\ref{prop:B-RD5} and~\ref{prop:B-RD18}.

Transport the first filtered subgroup of $N[p^2]$ to a cyclic
direct summand $C\subset M[p^2]$, with character
$\chi_p^b\epsilon$. If $b=0$, twist $M$ by $\epsilon^{-1}$;
the resulting height-two group has a rational point of order
$p^2$. If $b=1$, first form the BT quotient $M'=M/C$.
The image of $M[p^2]$ in $M'[p^2]$ is cyclic of order $p^2$
with character $\epsilon_0\epsilon^{-1}$; twisting by its
inverse again produces a rational point of exact order $p^2$.
The schematic closure of $C$ in a good model is finite flat,
and the quotient good BT group remains isogenous to the
original one. These quotients and finite twists therefore remain
potentially good, of dimension one and special slope $1/2$.

Lemma~\ref{lem:B-RD49-inertia} gives a good formal model over
an extension of ramification degree $m\mid p^2-1$.
Its absolute ramification is $a=(p-1)m$, which satisfies
(B-RD49.1). The rational torsion point extends by properness of
the finite flat torsion scheme, and reduces to the identity
because its special fibre is connected. It is consequently a
formal point forbidden by Lemma~\ref{lem:B-RD49-point}.
\end{proof}

\begin{theorem}[B-RD49: a one-step least-index interval in tame height two]
\label{thm:B-RD49}
\status{Proved; complete root and obstruction/crystal proof, original-source, and exact-code audits passed.}
For every prime $p\geq5$, in the setting of B.4 with
$e=p-1$, $\mu_p\subset K$, and height two,
\[
                         1\leq\nu_{\min}(p,p-1,2)\leq2.
                                                        \tag{B-RD49.9}
\]
More precisely, $p^3$-torsion determines the requested character
for every group satisfying the manuscript monodromy condition.
Any failure at $p^2$ must involve two whole potentially
supersingular groups. After completed maximal unramified base
change their finite determinant characters agree, and their
joint faithful tame inertia has degree
\[
                            m\mid p^2-1,\qquad m>p.    \tag{B-RD49.10}
\]
No existence of one of these remaining pairs is asserted.
\end{theorem}
\begin{proof}
The lower bound is Proposition~\ref{prop:B-RD18} with $n=1$.
Suppose two allowed groups have isomorphic $p^2$-torsion and
pass to the completed maximal unramified base. For either
monodromy filtration let $d\in\{0,1,2\}$ be the sum of the
dimensions of its good graded pieces. The generic determinant
on inertia is $\chi_p^d\epsilon$, with $\epsilon$ a finite
$\mu_{p-1}$-valued character. The image of $\chi_p$ modulo
$p^2$ has order $p$: the field contains $\mu_p$ but cannot
contain $\mu_{p^2}$, since that would require absolute
ramification divisible by $p(p-1)$. Equality of the two
determinants modulo $p^2$ thus gives equality of $d$, because
$|d_M-d_N|\leq2<p$, and then equality of $\epsilon$.

For $d=0$, the special character is the Brauer lift of the
semisimplified generic mod-$p$ inertia representation;
additivity makes this true for any extensions between the
graded pieces. For $d=2$, Cartier duality reduces to $d=0$.
For $d=1$, a potentially ordinary height-two graded group has
a canonical connected--etale filtration, which descends to
$K$, so each group is either filtered by height-one groups or
is itself potentially supersingular. Two filtered groups have
finite graded characters determined by the semisimplified
mod-$p$ representation, since $\chi_p\equiv1\pmod p$ on
inertia. The mixed case is excluded by
Proposition~\ref{prop:B-RD49-mixed}.

It remains to compare two whole potentially supersingular
groups. Lemma~\ref{lem:B-RD49-inertia} supplies a common good
extension with $e_L=(p-1)m$, $m\mid p^2-1$. As
\[
                     e_L\leq(p-1)(p^2-1)<(p-1)p^2,
\]
\cite[Theorem~E(1), pp.663--664; Proposition~18 and the full
proof of Theorem~E, p.685]{Bondarko2006Generic}, with $u=2$
and $\ell=1$, extends an isomorphism of generic $p^3$-torsion
to an isomorphism of the special $p$-torsion models. The
extension is compatible with the prescribed generic map and
therefore with descent. The resulting equal mod-$p$ special
Dieudonne representations have the same prime-to-$p$ Brauer
character, hence the same characteristic-zero finite inertia
character. This proves the upper bound. If $m\leq p$, the
same theorem with $u=1$, $\ell=1$ works already from
$p^2$-torsion and excludes a counterexample. The finite
norm characters agree by the determinant argument above,
which proves every restriction in (B-RD49.10).
\end{proof}

\begin{remark}[Scope and exact arithmetic check]
This argument includes nontrivial finite reduced norms and does
not replace arbitrary formal groups by elliptic curves. The
least value within (B-RD49.9) remains open here. The coefficient
normalization and universal deformation assertions used above
are read from the original Hopkins--Gross paper; the direct PDF
download and screenshot endpoint failed, while its indexed
original full text was successfully opened and read.
The retained program \path{code/verify_tame_height_two_gate.py}
checks the exact coefficient identity, the cancellation
arithmetic and the tame degree and determinant restrictions.
It reconstructed the $p$-series through degree $p^2$ for
$p=2,3,5,7,11$, checked $136220$ cancellation configurations,
$6984$ joint tame-degree cases and $1320$ determinant cases.
The author run passed in $1.3878417579981033$ seconds; independent
root and obstruction/crystal runs passed in $0.6928875380035606$
and $0.7361230790047557$ seconds, respectively. Their runs preserved
the author output. No existence claim for the remaining pairs is
inferred from these finite checks.
\end{remark}
